% ============================================================
% Chapter 4: Testing Methodology
% ============================================================
\chapter{Testing Methodology}
\label{chap:testing}

\section{Derivation of the Verification Strategy}
\label{sec:derivation_of_strategy}

The two-layer verification strategy of this project combines established models from functional safety and software engineering; it is not a new invention. Its structure is tailored to the challenges of a retrofit project, where a legacy function must be verified against modern standards.

The approach is based on the V-model, a lifecycle model commonly used for safety-critical development. The V-model posits that each phase of specification and design is mirrored by a corresponding verification or validation phase. This work maps the complementary principles of white-box and black-box testing along the V-model onto a hardware and firmware retrofit.

The methodology maps directly to the right side of this V-model (\Cref{fig:v_model_mapping}):
\begin{itemize}
	\item Layer~1, \emph{Component and Interface Tests}, corresponds to the lower levels of verification in the V-model: the unit and integration tests. This layer is meant to show that the individual design artefacts (\gls{vhdl} modules, hardware interfaces) are correctly implemented against the datasheets of the parts they connect to and the functions identified by reverse engineering; no requirement list exists for them (\Cref{sec:verification_artefacts}). This constitutes a classic white-box verification approach.
	\item Layer~2, \emph{System-Level Tests}, corresponds to the highest level of validation in the V-model: the system and acceptance tests. Here, the fully integrated system is to be validated as a black box against the original functional system requirements (the behaviour of the legacy system).
\end{itemize}

For planning, executing and documenting these tests, the objective was to follow the verification and validation processes of~\gls{ieee1012}~\cite{ieee1012_2024} (\Cref{sec:objectives_scope}); no~\gls{vv} plan in the sense of that standard was written, and no integrity level was assigned (\Cref{sec:safety_compliance}). This report combines these models into an application-specific strategy, in which new technology (the~\gls{fpga}) is verified and then validated against the behaviour of a legacy system.

Demonstrating the functional equivalence of the retrofitted anti-skid test-diagnosis module to its legacy counterpart~\cite{rexroth1992wsp}, with a view to a later assessment against contemporary safety standards, necessitated a structured and multi-faceted verification and validation strategy. With the verification and validation requirements of standards such as~\gls{iec61508}~\cite{iec61508_2010} and~\gls{ieee1012}~\cite{ieee1012_2024} in view, a two-layer testing methodology was defined for this safety-related retrofit; it has not been assessed against these requirements. It addresses the functions identified by reverse engineering the legacy system~\cite{rexroth1992wsp, ayoub2014, leyer2023}; for the modification itself it follows the split of \Cref{sec:en50716_modifications}, under which \gls{en50716}:2023~\cite{en50716_2023} bears only on the software side, the reused program and the \gls{stm32} firmware.

Central to this approach is the two-layer strategy (Layer~1, \textit{Component and Interface Tests}; Layer~2, \textit{System-Level Tests}), designed for targeted verification and progressive integration, so that issues can be isolated layer by layer. It combines detailed component-level verification (Layer~1) with system integration testing (Layer~2) on the way towards demonstrating equivalence. In the follow-on work, testing the rebuilt unit on the Layer~2 fixture exposed the symptoms of the most consequential defect, whose cause was then located by debugging (\Cref{sec:integration_defects}).

% ------------------------------------------------------------
\section{Layer 1: Component and Interface Tests}
\label{sec:fpga_interface_testing}

This foundational layer focused on verifying the internal logic of the~\gls{fpga} (\Cref{sec:fpga_core_details}) and its immediate electrical interfaces using deterministic, white-box methods. The primary goals were to confirm logical correctness, validate timing accuracy against expected behaviour derived from~\gls{mcs48} documentation and legacy system analysis~\cite{rexroth1992wsp}, and ensure basic hardware integrity before system integration. This approach is meant to follow, by analogy, the principles of~\gls{en50129} for safety-related electronic hardware, which here covers the \gls{hdl}/\gls{fpga} implementation (\Cref{sec:standards_detail}).

\begin{description}
	\item[VHDL Simulation and Testbenches:] Pre-synthesis simulations using the simulator bundled with the Microsemi Libero SoC toolchain (ModelSim ME, the Mentor/Siemens ModelSim Microsemi Edition) exercise the~\gls{fpga} design prior to hardware deployment. The~\gls{vhdl} testbenches, such as \texttt{tb\_testDiagnose.vhd}~\cite{leyer2023}, provide controlled and repeatable stimuli to the~\gls{dut}, that is, the top-level~\gls{fpga} design and its sub-modules. Key functions to be verified are: the logical behaviour of the emulated~\gls{mcs48} core by running test instruction sequences; memory interface timing (setup, hold and access times) of the external Flash~\gls{rom} (\texttt{S29AL016J}) and~\gls{nvram} (\texttt{CY14B104NA}) against their datasheets; correct operation of the emulated peripheral logic (7-segment display latching, button input and interrupt logic, frequency sweep state machine, \gls{spi}~master protocol); and the watchdog interaction logic (servicing of the watchdog via the chip-select line, response to \texttt{FPGA\_8048\_!RESET}). The expected behaviour is checked by manual waveform analysis: Leyer's digital appendix holds 17 testbench files of 2023 in 14 project folders, and none of them contains an assertion; not counted are two archives with files of 2020 in the same folder (\texttt{Achskarte.zip}, \texttt{RAMtest.zip}), and a testbench in one of them contains assertions. The archived simulation of the complete design completed without errors although the \gls{nvram} enable defect (B1, \Cref{tab:bug_taxonomy}) was present: the model of the \gls{nvram} in its testbench (\texttt{tb\_testDiagnose.vhd}) ignores the byte-low enable, so the simulation could not show B1 (\Cref{subsec:fpga_interface_results}). Simulation can therefore give no more than baseline confidence in the logic design and does not replace tests on the hardware.
	\item[JTAG and Pin Checks:] Reading the device identification (\texttt{A3P1000}) and programming and verifying the bitstream over~\gls{jtag} showed that the reused, desoldered~\gls{fpga} was alive; pin connectivity was checked with a test design that toggles neighbouring pins in opposite phase, measured with a multimeter~\cite{leyer2023}. Boundary scan (\gls{ieee1149}) was used to program the Flash~\gls{rom} through the~\gls{fpga} with \texttt{urjtag}; when this first failed, an investigation with a logic analyser and test designs ended in a visual inspection that found two cold solder joints, after which the Flash was programmed and verified~\cite[Sec.~4.3]{leyer2023}. These checks were made on the Rev~0.0 prototype in 2023 (\Cref{subsec:fpga_interface_results}).
	\item[Hardware Timing Measurements:] An oscilloscope is used to check critical timing parameters that are sensitive to the physical implementation: the watchdog reaction and the frequency sweep profile from \SI{850}{\hertz} to \SI{1550}{\hertz}. These measurements bridge simulated logical correctness and physical behaviour, because they check the actual RC networks, crystal tolerances and signal integrity against the design expectations. Leyer carried them out on the inherited Rev~0.0 prototype, comparing the sweep with that of the original board~\cite[Sec.~4.10]{leyer2023}; the watchdog of the original board could not be observed~\cite[Sec.~4.6]{leyer2023}. The archived captures are his (\Cref{fig:watchdog_reaction,fig:sweep_comparison}). On the reworked assembly, no such measurement is documented (\Cref{subsec:fpga_interface_results}).
\end{description}
This deterministic layer is intended to establish confidence in the basic functional correctness and timing behaviour of the~\gls{fpga} logic and its direct interfaces before the integrated system testing described next. Its archived simulations and measurements date from 2023, before the rebuild, and the simulations could not detect defect~B1 (\Cref{subsec:fpga_interface_results}).

% ------------------------------------------------------------
\section{Layer 2: System-Level Tests}
\label{sec:system_scenario_testing}

\begin{figure}[htbp]
	\centering
	% Layer 2 test infrastructure, kept as diagram source. Derived from the
% corresponding figure of the published paper. Drawn at its natural size:
% the gaps between the boxes are wider than every edge label, and the
% pass/fail log hangs off the host PC, whose application writes it.
\begin{tikzpicture}[
    blk/.style={draw, thick, rounded corners=3pt, minimum height=1.1cm,
                align=center, font=\footnotesize, text width=2.35cm},
    arr/.style={-{Stealth[length=2.5mm]}, thick},
    barr/.style={{Stealth[length=2.5mm]}-{Stealth[length=2.5mm]}, thick},
    lbl/.style={font=\scriptsize, text=black!70, align=center, inner sep=1.5pt},
    doc/.style={draw, thick, rounded corners=1pt, align=center,
                font=\scriptsize, inner sep=3pt, red!70!black, text=black},
  ]
  \node[blk, fill=blue!10]   (pc)  at (0,0)    {\textbf{Host PC}\\[-1pt]{\scriptsize \mbox{Python / YAML}}};
  \node[blk, fill=yellow!12] (nuc) at (4.4,0)  {\textbf{STM32}\\[-1pt]{\scriptsize Nucleo F401RE}};
  \node[blk, fill=orange!12] (fix) at (8.8,0)  {\textbf{Test Fixture}\\[-1pt]{\scriptsize Testkartentester}};
  \node[blk, fill=green!10]  (dut) at (13.2,0) {\textbf{DUT}\\[-1pt]{\scriptsize Retrofitted Board}};
  \draw[barr] (pc.east)  -- node[above, lbl] {USB/\\Serial}  (nuc.west);
  \draw[barr] (nuc.east) -- node[above, lbl] {GPIO}          (fix.west);
  \draw[barr] (fix.east) -- node[above, lbl] {Connector\\X4} (dut.west);
  \node[doc, below=0.6cm of pc] (log) {Pass/Fail log};
  \draw[arr, densely dashed, red!70!black] (pc.south) -- (log.north);
\end{tikzpicture}

	\caption[Layer 2 Test Infrastructure]{Layer~2 test infrastructure. The host PC orchestrates YAML-defined scenarios via an \gls{stm32} Nucleo board, which is connected to the dedicated test fixture~\cite{dolata2023} by general-purpose input/output (\acrshort{gpio}) lines; the fixture stimulates and monitors the device under test (\acrshort{dut}) through the system connector \texttt{X4}.}
	\label{fig:layer2_setup}
\end{figure}

Building upon component verification, this layer is meant to validate the integrated system (the Base Board with the programmed~\gls{fpga}, the Upper Board, the~\gls{stm32} Diagnostics Board and the LED Display Board) under operational scenarios using a dedicated hardware test fixture (\Cref{fig:layer2_setup}). It addresses a fundamental limitation of Layer~1: simulation and hardware measurements verify individual components against their specifications, but cannot confirm that the complete system reproduces the legacy behaviour that operators and maintenance personnel rely on. This layer adopts a black-box perspective at board level. On this setup the follow-on team observed the symptoms of the \gls{nvram} enable defect at the board boundary (\Cref{sec:integration_defects}).

\begin{description}
	\item[Dedicated Test Fixture (Testkartentester):] Testing relied on the hardware test fixture developed by Dolata et al.~\cite{dolata2023}. The fixture provides the physical connection to the retrofitted board assembly through the system connector \texttt{X4} (48-pin DIN~41612, type~F) and interfaces with a PC via an intermediate~\gls{stm32} Nucleo board (\texttt{STM32F401RE}, see~\cite[Sec.~6.2]{dolata2023}). It allows the PC-based test software to stimulate the board's inputs (e.g., simulating button presses or signal conditions corresponding to legacy faults using pins such as \texttt{18z}, \texttt{24d}) and to monitor its outputs.
	\item[Four-Phase Test Cycle:] Each test scenario follows the same cycle: (1)~stimulus application via pin control, (2)~a settling delay (typically \SI{3}{\second}, matching the legacy button timing), (3)~assertion of the expected output, and (4)~state cleanup via the Lösch (clear) button.
	\item[YAML-Based Scenario Definition:] To ensure structured and repeatable system-level testing, test procedures are defined in the human-readable~\gls{yaml} format of the fixture~\cite{dolata2023} (\Cref{sec:app_yaml_test_cases}). Each test case specifies the input stimuli to be applied via the fixture (\texttt{pin\_sets}), the expected system state (primarily the code shown on the 7-segment displays, \texttt{expected\_err\_code}), the expected stored fault sequence (\texttt{expected\_err\_seq}, read out via the STW button), and prompts for manual tester interaction or verification where necessary (e.g., confirming the displayed code or pressing the \enquote{Lösch} button, flag \texttt{del: True}). Each case is named after the fault it simulates and expects the codes documented for that fault in~\cite{rexroth1992wsp}, which links the cases to the fault codes by name; the check of the self-test was still outstanding in the fixture of 2023~\cite[Sec.~6.2]{dolata2023}. For example, simulating a \enquote{Geber~3} fault (sensor fault on axle card~3) sets \verb|pin_sets: {18z: 1, 24d: 0}| and expects the code \texttt{31} on the display and the sequence \texttt{03~31~09}~\cite{dolata2023}; for the display, the manual and the table of the fixture report give \texttt{03} (\Cref{tab:requirement_basis}, note~b), so the \texttt{31} is the fixture's expectation, not the legacy behaviour.
	\item[Interactive Test Execution:] The fixture's Python application (\texttt{main.py}, with modules such as \texttt{run\_tests.py} and \texttt{load\_tests.py})~\cite{dolata2023} runs on a PC and orchestrates the test execution based on the loaded~\gls{yaml} scenarios. It communicates with the test fixture via the Nucleo board (using \texttt{pyserial}) to apply the defined stimuli, pauses execution in a text-based user interface (\texttt{gui/testing.py}) to prompt the tester for manual actions or observations (e.g., \enquote{Press STW Button}, \enquote{Verify Display shows 31}), and logs the pass/fail outcome with timestamps (\texttt{logger/test\_logger.py}, see~\Cref{sec:app_system_log_excerpts}). In this semi-automated procedure the tester checks outputs such as the displays by eye and confirms each step.
	\item[Coverage:] The test suite is designed to cover (1)~complete fault memory cycles (trigger $\rightarrow$ display $\rightarrow$ store $\rightarrow$ recall $\rightarrow$ clear) for the fault codes of the four axle cards in the legacy diagnosis table~\cite[Blatt~20/25]{rexroth1992wsp}, (2)~self-test activation and frequency sweep, (3)~button interaction equivalence, (4)~power-cycle stability, and (5)~diagnostic~\gls{spi} transmission to the~\gls{stm32} logger. In detail:
	      \begin{itemize}
		      \item \textit{Complete Fault Memory Cycle:} Triggering specific fault conditions via input pins (e.g., the Geber~3 fault via pins \texttt{18z}/\texttt{24d}), verifying the code on the display (the fixture expects \texttt{31}, \Cref{tab:requirement_basis} gives \texttt{03}), simulating the \enquote{STW} button press to read the stored error sequence (\texttt{03~31~09}), verifying storage and recall across power cycles (the fixture switches the supply of the board under test through relays~\cite{dolata2023}), and confirming the clearing of faults with the simulated \enquote{Lösch} button (resulting in the \enquote{99} display).
		      \item \textit{Self-Test Activation and Execution:} Simulating the \enquote{Test} button press to initiate the legacy self-test sequence, verifying the sequence of codes displayed (\Cref{subsec:requirement_basis}), and confirming the activation of the frequency sweep output during the appropriate phase (timing measured in Layer~1).
		      \item \textit{User Interface Equivalence:} Verifying that the system's response to simulated button presses (Test, STW, Lösch, Tür/V) and the resulting display states correspond to the behaviour described in the original system manual~\cite{rexroth1992wsp}.
		      \item \textit{Integrated System Stability:} Running sequences of tests and power cycles to check for unexpected hangs or failures in the integrated system.
		      \item \textit{Diagnostics Subsystem Interaction:} Verifying that diagnostic frames are transmitted from the~\gls{fpga} to the~\gls{stm32} board via the isolated~\gls{spi} link (\Cref{sec:stm32_details}) during fault scenarios and logged appropriately (\Cref{subsec:system_scenario_results}).
	      \end{itemize}
\end{description}
These fixture-based tests are central to demonstrating functional equivalence under relevant operational conditions, providing system-level evidence complementing the component-level verification of Layer~1; how far they were run on the reworked assembly is described in~\Cref{subsec:system_scenario_results}.

\subsection{Requirement Basis}
\label{subsec:requirement_basis}

The requirements against which the module is checked are the behaviour of the legacy module as its manual and the fixture report describe it; \Cref{tab:requirement_basis} collects the diagnosis codes. As this report uses the term, functional equivalence means that, for the same inputs at the system connector and the same button actions, the retrofitted module shows the codes and read-out sequences of this table and runs the legacy self-test, whose frequency sweep is compared with that of the original board. No tolerances are defined, so the end frequencies and the duration of the sweep, the rhythm of the display and the duration of the self-test cannot be judged against the legacy behaviour. The sources give nominal values only: \SI{850}{\hertz} and \SI{1550}{\hertz} as the end frequencies of the original's sweep~\cite[Sec.~2.9]{leyer2023}, a change from \SI{1540}{\hertz} to \SI{920}{\hertz} lasting about \SI{180}{\milli\second} on the original test card~\cite[Sec.~5.1.5, 5.2]{dolata2023}, three seconds per display in the read-out and for the lamp test~\cite[Blatt~14/25]{rexroth1992wsp}, and about one minute for the test program~\cite[Blatt~14/25, 18/25]{rexroth1992wsp}, which the fixture report gives as almost \SI{70}{\second}~\cite[Sec.~5.1]{dolata2023}. The table is a requirement basis, not a test specification (\Cref{subsec:equivalence_discussion}).

\begin{table}[htbp]
	\centering
	\caption[Legacy Diagnosis Codes as Requirement Basis]{Legacy diagnosis codes as requirement basis; X is the number of the axle card (1 to 4) or, for the computer cards, of the card (1 to 5).}
	\label{tab:requirement_basis}
	\small
	\begin{tabular}{@{}>{\raggedright\arraybackslash}p{0.33\textwidth} >{\raggedright\arraybackslash}p{0.12\textwidth} l >{\raggedright\arraybackslash}p{0.13\textwidth} >{\raggedright\arraybackslash}p{0.2\textwidth}@{}}
		\toprule
		Fault or state$^{a}$ & Code & Display & Read-out (STW) & Rebuilt unit$^{c}$ \\
		\midrule
		Sensor: short circuit or wire break & 11, 21, 31, 41 & 0X$^{b}$ & 0X, X1, 09 & In the suite, not recorded \\
		Solenoid valves: short circuit & 13, 23, 33, 43 & 0X & 0X, X3 & In the suite, not recorded \\
		Solenoid valves: wire break & 14, 24, 34, 44 & 0X & 0X, X4 & In the suite, not recorded \\
		Bogie monitoring time & 10, 20, 30, 40 & 0X & 0X, X0, 05 & In the suite, not recorded \\
		Computer card (05: comparator card) & 15, 25, 35, 45, 55 & 0X & 0X, X5 & In the suite, not recorded \\
		Bogie monitoring time in the self-test & -- & 95 & 95, 0X, 05 & Not documented \\
		All monitoring times in the self-test & -- & 95 & 95, 05, 06 & Not documented \\
		\midrule
		No fault & 99 & & & Not documented \\
		Sporadic fault (no longer active) & 95 & & & In the suite, not recorded \\
		Lamp test & 88 & & & Not documented \\
		Test run (self-test) & 89 & & & Not documented \\
		\bottomrule
	\end{tabular}

	\medskip
	\begin{minipage}{0.95\textwidth}
		\footnotesize
		$^{a}$Codes, displays and read-out sequences after the diagnosis table of the manual~\cite[Blatt~20/25]{rexroth1992wsp}. The fixture report gives the same values~\cite[Tab.~2]{dolata2023}, except that for the two self-test faults it gives the read-out without the leading 95.\\
		$^{b}$For the sensor fault of axle card~3, the fixture's test case expects 31 on the display while the fault is applied~\cite[Sec.~9.3, Listing~3]{dolata2023}, where both tables give 03; the sources do not explain the difference.\\
		$^{c}$The follow-on team ran the fixture's full suite and 19 selected cases on the rebuilt unit and saw the expected codes, without recording which codes the cases contained~\cite[p.~15]{savchuk2025}. The self-test was not among the fixture's cases of 2023~\cite[Sec.~6.2]{dolata2023}; on the fixture firmware used in 2025 see \Cref{subsec:system_scenario_results}.
	\end{minipage}
\end{table}

The self-test runs after the supply is switched on, or at standstill when the Test button is pressed with the braking signal present: the display changes from 99 to 88 (lamp test) and to 89 (test run), the test program runs for about one minute and makes the anti-skid control valves act one after another, and afterwards 99 must appear again if all systems are in order, otherwise a fault code~\cite[Blatt~14/25, 18/25, 19/25]{rexroth1992wsp}. Leyer expected 85 after the 88~\cite[Sec.~5.1.2]{leyer2023}, where the manual gives 89; the sources do not explain the difference. When STW is pressed, the legacy module shows 88 for three seconds, then the stored information in a rhythm of three seconds, and stops at the first display; Lösch clears all stored information~\cite[Blatt~14/25]{rexroth1992wsp}.

% ------------------------------------------------------------
\section[Validation Context and Complementary Approaches]{Validation Context and Complementary\\Approaches}
\label{sec:validation_context}

The primary two-layer testing strategy was situated within a broader validation context, acknowledging both the project's scope and the requirements for safety-critical systems.

\begin{description}
	\item[Relation to \gls{hil} Testing:] A full dynamic vehicle model interacting with the hardware in real time was beyond the project's scope and resources. The method borrows from~\gls{hil} practice~\cite{hil_rail_suspension_2023, zhou_wu_2020}: it combines component-level simulation (Layer~1) with scripted functional tests of the real hardware on a fixture that applies simulated fault inputs (Layer~2). The tester reads the displays and marks each case as passed or failed by hand~\cite[Sec.~9.2]{dolata2023}, and no plant model closes the loop. The fixture could be a starting point for a~\gls{hil} test rig (\Cref{sec:future_work}).
	\item[Fault Handling Verification:] The focus of fault testing in Layer~2 was on verifying the system's detection of and response to the known legacy error conditions (as defined in~\cite[Blatt~20/25]{rexroth1992wsp} and triggered via the fixture inputs defined in the \gls{yaml} test cases). This aligns directly with the primary goal of demonstrating functional equivalence. Exhaustive low-level hardware fault injection (e.g., simulating stuck-at faults or bit flips within the~\gls{fpga} using techniques surveyed in~\cite{ziade2004faultinjection}) was not performed, as the main objective was replicating existing behaviour rather than proving resilience against arbitrary internal hardware faults beyond the scope of the original design's fault detection mechanisms.
	\item[Formal Verification Scope:] Verification relied primarily on simulation (Layer~1) and testing (Layers~1 and~2). Formal methods, such as model checking or formal equivalence checking~\cite{formal_verification_rail_2023, mneimneh2005sequential}, were not applied to the~\gls{hdl} or software components due to project scope and complexity. The~\gls{vhdl} testbenches were meant to check logical correctness and to reveal design flaws early in the process; the archived testbenches of 2023 contain no assertions (\Cref{sec:fpga_interface_testing}), and no code review is documented.
	\item[Test Coverage Considerations:] The verification strategy aimed for \textit{functional coverage} based on the legacy system's requirements and documented behaviour~\cite{rexroth1992wsp}. The archived Layer~1 testbenches of 2023 exercise the complete design and the test designs for individual functions, such as the sweep and the watchdog reset test; none of them contains assertions, and their results were checked by eye (\Cref{sec:fpga_interface_testing}). Layer~2 scenario testing aimed for requirements-based coverage and was intended to check the documented fault codes of the axle cards, the self-test sequence and the user interactions. Quantitative code coverage metrics (e.g., statement and branch coverage for C, structural coverage for \gls{vhdl}) were not systematically measured due to tool limitations and scope; the methodology instead aimed at confidence through targeted functional and scenario-based validation. Test cases derived from the control-flow structure, as in structured testing~\cite{watson1996}, were not used; together with structural coverage measurement they are an option for later work (\Cref{sec:future_work}).
	\item[Environmental Qualification Approach:] Compliance with the standards for rolling-stock equipment named in \Cref{sec:standards_detail} (\gls{en50155}:2021~\cite{en50155_2021}, \gls{en61373}:2010~\cite{en61373_2010} and the \gls{en50121} series~\cite{en50121_3_2_2016}, cited by number and title) has not been tested and is not claimed; the design choices made following established principles (\Cref{sec:design_principles}) are the only step towards it. A formal demonstration of compliance would need tests of the finalised hardware against these standards, for example by an accredited laboratory; they were outside the scope of this research phase.
\end{description}
These complementary approaches define the boundaries and context of the verification performed, highlighting the focus on demonstrating functional equivalence through simulation and fixture-based testing, while acknowledging the areas requiring further work (such as accredited environmental testing) for full certification.

% ------------------------------------------------------------
\section{Verification Artefacts and Process Control}
\label{sec:verification_artefacts}

Generating auditable evidence suitable for supporting a safety case according to standards such as~\gls{iec61508}~\cite{iec61508_2010}, \gls{en50129}~\cite{en50129_2018}, \gls{en50716}~\cite{en50716_2023}, and~\gls{ieee1012}~\cite{ieee1012_2024} was a key process objective throughout the verification activities.

\begin{description}
	\item[Test Documentation and Logs:] The core artefacts generated and managed include:
	      \begin{itemize}
		      \item \gls{vhdl} testbenches: stimulus files for the Layer~1 simulations, without assertions (e.g., \texttt{tb\_testDiagnose.vhd}; \Cref{sec:fpga_interface_testing}); \Cref{app:test_cases} shows the \gls{nvram} control of the top-level design.
		      \item Measurement captures: the archived oscilloscope captures of the watchdog test and of the frequency sweep, taken with the Layer~1 method on the inherited prototype~\cite{leyer2023} (\Cref{fig:watchdog_reaction,fig:sweep_comparison}).
		      \item \gls{yaml} test definitions (\Cref{app:test_cases}): scenarios for Layer~2 system testing in the format of the fixture~\cite{dolata2023}; the definitions used in the follow-on runs are not part of the archived material.
		      \item Fixture test logs (\Cref{app:test_logs}): the fixture application records executed Layer~2 tests with tester inputs and pass/fail results (format in~\Cref{sec:app_system_log_excerpts}); no log of the follow-on runs is part of the archived material.
		      \item Hardware measurement records: none from the reworked assembly (\Cref{tab:results_summary}).
	      \end{itemize}
	\item[Traceability:] No requirement list exists for the module, and no traceability matrix was generated. The Layer~1 testbenches and test designs were derived from the functions identified by reverse engineering, such as the frequency sweep, the watchdog, the \gls{nvram} and the Flash interface~\cite[Sec.~2]{leyer2023}; they are stimulus files without assertions (\Cref{sec:fpga_interface_testing}), and no module requirement is documented for them to trace to. The Layer~2 YAML scenarios of the fixture are linked to the legacy behaviour documented in~\cite{rexroth1992wsp} by the name of the fault case: each case names a fault of the diagnosis table and gives its expected codes, and the STW and Lösch functions take part in the cases~\cite[Sec.~9.3, Listing~3]{dolata2023}. The self-test was not among the cases of 2023 (\Cref{chap:safety_checklist}, item~CL-D1). The verification results (\Cref{chap:results}) link these tests back to the design elements (\Cref{chap:architecture}). For formal certification, requirements would have to be specified and linked to the tests in explicit traceability matrices; the checklist in~\Cref{chap:safety_checklist} gives a high-level mapping as a step towards this.
	\item[Version Control:] The follow-on work kept its sources in Git~\cite[p.~7]{savchuk2025}; the test definitions and logs of the Layer~2 runs are not part of the archived material. A certification campaign would have to place all design files, test artefacts, firmware and documentation under configuration management, supporting reproducibility and change tracking.
	\item[Tool Qualification Considerations:] Standard, widely used tools (Microsemi Libero SoC including ModelSim ME, STM32CubeIDE, Python with libraries such as \texttt{pyserial}) were employed. For the software tools, STM32CubeIDE and the Python software of the fixture, a formal certification according to~\gls{iec61508} or~\gls{en50716} would call for a tool qualification assessment: \mbox{\gls{iec61508}-3} provides requirements for support tools such as development and design tools, language translators, and testing and debugging tools~\cite{iec61508_3_2010}; each tool would be classified (e.g., T1, T2, T3 by potential impact) and its use justified, potentially through vendor qualification data, validation suites, or evidence of operational experience. The \gls{fpga} tools, Libero SoC with ModelSim ME, serve the \gls{hdl}, which is treated as hardware here; \gls{en50716} does not cover them, and \Cref{subsec:fpga_certification} says what is open about them. No tool of this project has been assessed.
\end{description}

\section{Verification of Modifications}
\label{sec:en50716_modifications}

A retrofit is a modification of an existing system. In this one the controller hardware is replaced and the program is kept: \gls{en50716} applies only to software and its interaction with the system~\cite{en50716_2023}, so it bears on the reused program and the~\gls{stm32} firmware, while the change of hardware falls to the hardware standards (\Cref{sec:standards_detail}). The verification methodology addresses the four elements of change management below; it has not been assessed against the requirements of these standards.

\begin{description}
	\item[Impact Analysis:] No impact analysis is documented. The scope of the verification follows from the premise of the design: the modification is confined to the hardware of the test-diagnosis board, while the original program of the emulated controller is kept unchanged (\Cref{sec:fpga_core_details}). Two of the hardware changes alter the behaviour of the module: a~\gls{bod} circuit was added to the reset circuit kept from the original board~\cite[Sec.~2]{leyer2023}, and the~\gls{nvram} was replaced by a part that needs no STORE and RECALL signals~\cite[Sec.~2.5]{leyer2023}. On the original test card, the STORE and RECALL inputs of the system connector (08z, 08d) trigger the storing and the recalling of the fault memory~\cite[Sec.~5.1.4]{dolata2023}; in the~\gls{fpga} design of 2023 they feed only the logger (\texttt{atop.vhd}), so they no longer act on the fault memory. Whether a reset in service leads into the automatic check is open (\Cref{sec:vehicle_outputs}).
	\item[Verification of the Change:] The two-layer methodology is designed as the process for verifying the modification: Layer~1 for the correct implementation of the new hardware and~\gls{hdl} components at a technical level, and Layer~2 for the system behaviour that the change has to keep.
	\item[Regression Testing:] The Layer~2 scenario tests are meant to serve as a regression test suite for the documented legacy functions. Run and logged after each change, they would re-validate the known legacy functions (fault memory, self-tests) and so show that the modification has not introduced unintended adverse effects on existing, functional parts of the system; in this project neither a log of the runs on the rebuilt unit nor a run of the self-test on that unit is documented (\Cref{tab:evidence_state}).
	\item[Configuration Management:] As detailed in \Cref{sec:verification_artefacts}, configuration management of the artefacts (hardware design, \gls{hdl}, software, test cases, documentation) supports the traceability and reproducibility of changes; in this project it is incomplete, because the Layer~2 test definitions and logs were not archived.
\end{description}

The two layers and the process controls described above are meant to provide an auditable evidence base; the evidence archived in this project and its gaps are discussed in the following chapter.
